\chapter{夸克质量方案的讨论}
\label{chap:mass_scheme}
\section{引言}
这里写从拉格朗日量出发的巴拉巴拉，引入下面的极质量介绍，当然先简单写一些轻夸克质量方案的介绍，最后介绍对于重夸克质量的naive方案，pole mass方案的简单介绍。

本章的核心问题不是简单比较不同质量方案的定义，而是从 threshold $Q\bar Q$ 体系与静态总能量的物理图像出发，理解 pole mass 的局限、1S/PS 等阈值短距质量的构造逻辑，以及 $\overline{\mathrm{MS}}$ 质量在高能短距过程中的自然性。

\section{极质量方案}
极质量方案基于重夸克自能过程(图\ref{fig:bubble-renormalon})，定义为重夸克在传播子极点处的质量。该定义基于在壳重整化方案，要求修正后的重夸克质量等于物理质量，修正后的电荷等于物理电荷，传播子在极点处的留数等于1，并且在微扰论的逐阶都是well-defined。事实上，夸克质量并不同于电子质量，由于QCD中的夸克禁闭现象，夸克并非渐近态，且QCD的S矩阵在$p^2=m^2$处并不存在夸克极点。这种“非物理性”导致QCD中的极质量不可避免的带有$\mathcal{O}(\Lambda_{QCD})$量级的不确定性，这种不确定性表现为领先的红外重整子的模糊性问题。

\begin{figure}[htbp]
\centering
\begin{tikzpicture}[x=1cm,y=1cm,line cap=round,line join=round]

  \tikzset{
    hq/.style={line width=0.9pt},
    gluon/.style={
      draw,
      decorate,
      decoration={coil,aspect=0.62,segment length=3.0mm,amplitude=1.45mm},
      line width=0.45pt
    },
    dgluon/.style={
      draw,
      dashed,
      decorate,
      decoration={coil,aspect=0.62,segment length=3.0mm,amplitude=1.45mm},
      line width=0.45pt
    },
    bubble/.style={
      circle,
      draw,
      minimum size=0.82cm,
      inner sep=0pt,
      pattern=north east lines,
      pattern color=black,
      line width=0.45pt
    },
    momarrow/.style={
      postaction={
        decorate,
        decoration={
          markings,
          mark=at position #1 with {\arrow{Latex[length=2.6mm,width=1.8mm]}}
        }
      }
    },
    two momarrows/.style args={#1,#2}{
      postaction={
        decorate,
        decoration={
          markings,
          mark=at position #1 with {\arrow{Latex[length=2.6mm,width=1.8mm]}},
          mark=at position #2 with {\arrow{Latex[length=2.6mm,width=1.8mm]}}
        }
      }
    }
  }

  % 外部夸克线（左右各一个 p 箭头）
  \draw[hq,two momarrows={0.12,0.92}] (0,0) -- (8.4,0);
  \node at (1.05,-0.42) {$p$};
  \node at (7.35,-0.42) {$p$};

  % 左侧第一段胶子线（圈动量 k）
  \draw[gluon,momarrow=0.58]
    (1.78,0.00)
    .. controls (1.76,0.30) and (1.84,0.72) ..
    (2.07,1.02);
  \node at (1.35,0.78) {$k$};

  \node[bubble] (B1) at (2.18,1.28) {};

  % 中间虚线胶子
  \draw[dgluon]
    (2.49,1.55)
    .. controls (2.95,2.06) and (3.65,2.34) ..
    (4.30,2.28);

  \node[bubble] (B2) at (4.62,2.20) {};

  % 右侧上方胶子
  \draw[gluon]
    (4.93,1.98)
    .. controls (5.18,1.78) and (5.45,1.50) ..
    (5.67,1.24);

  \node[bubble] (B3) at (5.96,1.16) {};

  % 右侧胶子回到底部夸克线
  \draw[gluon]
    (6.25,0.92)
    .. controls (6.46,0.70) and (6.57,0.34) ..
    (6.52,0.00);

  % 内部夸克传播子（动量 p+k）
  \draw[hq,momarrow=0.63] (1.78,0) -- (6.52,0);
  \node at (4.15,0.34) {$p+k$};

  % 顶点
  \fill[blue] (1.78,0.00) circle (3.4pt);
  \fill[blue] (6.52,0.00) circle (3.4pt);

\end{tikzpicture}
\caption{夸克自能修正图中的链式泡沫图}
\label{fig:bubble-renormalon}
\end{figure}
仿照照量子电动力学中计算电子自能过程的标准技术，最终该图的贡献是
\begin{equation}
    -i \Sigma(\slashed{p})
=
(-i g_s)^2 C_F \mu^{2\epsilon}
\int \frac{d^d k}{(2\pi)^d}\,
\gamma_\mu\,
\frac{i(\slashed{p}+\slashed{k}+m)}{(p+k)^2-m^2+i0}\,
\gamma_\nu\,
\frac{-i g^{\mu\nu}}{k^2+i0}\,.
\label{eq:self-energy}
\end{equation}
以单粒子不可约费曼图为例，对外部动量为 $p$ 的夸克线，取圈动量为 $k$，则内部夸克传播子的动量为 $p+k$。两个夸克——胶子顶点分别贡献因子 $-ig_s\gamma^\mu T^a$，内部夸克传播子贡献
\[
\frac{i(\slashed p+\slashed k+m)}{(p+k)^2-m^2+i0},
\]
在 Feynman 规范下，内部胶子传播子贡献
\[
\frac{-i\,\delta^{ab}g_{\mu\nu}}{k^2+i0}.
\]
再乘上维数正规化下的圈积分测度 $\mu^{2\epsilon}\!\int d^d k/(2\pi)^d$，并将颜色指标收缩为$T^aT^a=C_F,$
便得到式\eqref{eq:self-energy}的具体形式。 
其中 $d=4-2\epsilon$，$\mu$ 是维数正规化中引入的 ’t Hooft 质量尺度\footnote{在维度正规化方案中引入质量尺度$\mu$这一现象同时存在于QED与QCD的单圈自能修正计算过程。原因在于积分维度的改变引起规范耦合常数携带量纲，$\mu$是量纲补偿因子，使得$g_0=\mu^{\epsilon}Z_g g$不携带量纲。$\mu$同时是重整化尺度(量纲为1)，\textbf{物理可观测量若考虑完整的微扰修正，则不应该对$\mu$存在依赖。换而言之，有限阶截断后会保留残余的重整化尺度依赖，这也被用来估计高阶修正不确定性的来源之一}。具体地，若考虑完整微扰修正的物理可观测量$\mathcal{O}$，其不存在对重整化尺度的依赖，因此$\mu\frac{d}{d \mu}\mathcal{O}=0$。对于有限阶微扰修正的截断，该方程并非严格为0，这个不为零的量对应被截断的更高阶修正。换而言之，将某个可观测量写成标准微扰展开形式$O_N(\mu)=\sum_{n=0}^N c_n\left(\ln \frac{\mu}{Q}\right) \alpha_s^{n+p}(\mu)$, $p$是该观测量的微扰展开从$\alpha_s$的第几阶出现，其精确结果不存在对重整化尺度的依赖。但对于有限阶修正，其满足$\mu \frac{d}{d \mu} O_N=\mathcal{O}\left(\alpha_s^{N+p+1}\right)$.该方程说明，对于有限阶截断的物理可观测量，其对重整化尺度的敏感性本身就是”下一阶乃至更高阶未算项“。因此对重整化尺度在某个合理区间里遍历取值得到物理可观测量的数值改变，自然可以被拿来作为 missing higher orders 修正的经验估计，也就是所谓的 \textbf{renormalization scale variation}}。

对于式\eqref{eq:self-energy},假设外部动量$p$接近于夸克质量$m$，即on-shell条件成立，讨论软贡献区域，即$k^\mu$远小于m，则此时的重夸克传播子具有eikonal近似形式
$$
\frac{1}{(p+k)^2-m^2+i 0} \rightarrow \frac{1}{2 m v \cdot k+i 0}
$$
在该近似下，式\eqref{eq:self-energy}的软贡献可以被重写为
$$
\Delta m_{\mathrm{soft}} \sim \int \frac{d^d k}{(2 \pi)^d} \frac{1}{k^2} \frac{1}{v \cdot k}\sim\int k^3 \frac{1}{k^3}dk \sim \int_0^\mu d k \text { (constant) }
$$
分子在$k\to 0$极限下没有奇异项，仅贡献常数或者k趋于0的高阶修正。因此这里仅讨论分母。红外区域给出的是IR的线性依赖，后面再讨论泡泡链的修正后，会有$\Delta m_{\mathrm{IR}}^{(n)} \sim \int_0^\mu d k \ln ^n\left(\frac{k^2}{\mu^2}\right)$的贡献。

以上是关于单胶子泡泡图的简单说明，对于多胶子链式泡泡图，首先关注泡泡之间的传播子连接，一般意义上的胶子传播子具有如下的形式
$$
D_{\mu \nu}^{(0)}(k)=\frac{-i}{k^2}\left(g_{\mu \nu}-(1-\xi) \frac{k_\mu k_\nu}{k^2}\right) .
$$
将泡泡的动力学信息编码为$\Pi$,则插入一个泡，会有$D^{(0)}\Pi D^{(0)}$，插入两个泡会有$D^{(0)}\Pi D^{(0)}\Pi D^{(0)}$，以此类推，最终链式泡泡图的贡献可以被重写为
$$
D=D^{(0)}+D^{(0)} \Pi D^{(0)}+D^{(0)} \Pi D^{(0)} \Pi D^{(0)}+\cdots
$$
形式上根据Dyeson级数求和可并展开为几何级数,
$$D=\frac{D^{(0)}}{1+\Pi(k^2)}=D^{(0)}\sum_{M=0}^{\infty}[-\Pi(k^2)]^N$$
对于泡泡$\Pi_{k^2}$,由规范对称性和横向极化要求，$\Pi_{k^2}$具有如下的结构
$$\Pi_{\mu \nu}(k)=\left(k_\mu k_\nu-k^2 g_{\mu \nu}\right) \Pi\left(k^2\right),$$
表明该无量纲标量函数$\Pi(k^2)$在无质量极限下(大$n_f$极限下)，圈积分唯一的尺度依赖是$k^2$。根据标准圈图计算结果，通过量纲分析，可知$\Pi(k^2)$具有如下的一般结构
\begin{equation}
\Pi\left(k^2\right)=\frac{\alpha_s}{4 \pi}\left[\text { divergence }+ \text { const }+c \ln \frac{-k^2}{\mu^2}\right] .
\end{equation}
该圈积分的贡献只与对数项和一个重整化方案依赖的常数项相关。最终该含有N个泡泡的图会给出$\ln^{N}(-k^2/\mu^2)$的项.

在散射振幅的水平，胶子传播子往往和两个有效顶点耦合在一起，形成一个有效的耦合常数，记为$\alpha_s(k^2)$，
\begin{equation}
  \alpha_s(k^2)=\frac{g_s^2}{1+\Pi(k^2)}\propto \frac{g_s^2}{1+\frac{\beta_0 \alpha_s}{4 \pi} \ln \frac{k^2}{\mu^2}}+\cdots,
  \label{eq: alphas_running}
\end{equation}

此外，$\alpha_s$随尺度$\mu$的变化由重整化群方程描述，$\mu \frac{d\alpha_s}{d\mu}=\beta$,单圈水平下该方程的解是
$\alpha_s(k^2)=\alpha_s(\mu^2)/(1+\beta_0 \alpha_s(\mu^2)/(4 \pi) \ln(k^2/\mu^2))$

泡泡链重求和产生的传播子修正与重整化群方程的一圈解呈现相同的对数结构，因此可以将其理解为耦合常数重整化群跑动的图像化体现\cite{Papavassiliou_1997}。

通过引入跑动耦合常数$\alpha_s(k)$，质量位移的软贡献可以重新写为：
\begin{equation}
  \Delta_{\mathrm{soft}} \sim K\int_{0}^{\mu} \alpha_s(k) dk, 
\end{equation}
其中K是可计算的常数，比如在$\overline{\text{MS}}$方案下，$K=C_{F} e^{5/6}/\pi$.
\subsection{Renormalon: 阶乘增长 ,Borel变换与模糊性}
事实上，pole mass在定义在极点处，微扰论意义上是逐阶良定义的，但依然存在软贡献区域存在线性的IR敏感性，表现为阶乘形式的增长。
将式\eqref{eq: alphas_running}按照几何级数在软贡献区域进行展开，有$$\alpha_s(k)=\alpha_s(\mu) \sum_{n=0}^{\infty}\left[\beta_0 \alpha_s(\mu) \ln \left(\mu^2 / k^2\right)\right]^n$$.质量位移计算结果可改写为
$$
\Delta m_{\mathrm{soft}} \sim K \sum_{n=0}^{\infty} \alpha_s^{n+1}(\mu) \beta_0^n \int_0^\mu d k \ln ^n\left(\mu^2 / k^2\right)
$$
该结果对应的特征积分是$I_n=\int_0^\mu d k \ln ^n\left(\mu^2 / k^2\right)$, 做如下的变量代换，有
$$
t=\ln (\mu / k) \Rightarrow k=\mu e^{-t}, d k=-\mu e^{-t} d t
$$
$I_n$可以被重写为
\begin{equation}
  \boxed{
I_n=\mu \int_0^{\infty} d t e^{-t}(2 t)^n=\mu 2^n \Gamma(n+1)=\mu 2^n n!}
\label{eq: chatacter_Integral}
\end{equation}
因此，质量位移的软贡献具有阶乘增长的行为\cite{Bigi:1994em}。以$\overline{\text{MS}}$质量方案为例，最终得到的质量位移结果具有如下的结\cite{Beneke:2021lkq}
\begin{equation}
m_{\mathrm{pole}}-m_{\overline{\mathrm{MS}}}(\mu)=\frac{C_F e^{5 / 6}}{\pi} \mu \sum_{n \geq 0}\left(2 \beta_0\right)^n n!\alpha_s^{n+1}(\mu),
\end{equation}

Borel变换将阶乘增长的级数转换为Borel平面的奇点结构（如图\ref{fig:borel_plane_adler_pretty}所示），从Borel平面可直接读出渐进级数为什么发散，来自什么物理区域的贡献，以及对应什么量级的不确定性，这也是引入Borel变换的动机。
\begin{figure}[htbp]
\centering
\begin{tikzpicture}[scale=1.05, >=stealth]

    % axes
    \draw[->, thick] (-0.8,0) -- (8.2,0) node[right] {$\mathrm{Re}\,u$};
    \draw[->, thick] (0,-2.2) -- (0,2.6) node[above] {$\mathrm{Im}\,u$};

    % branch points and cuts
    \foreach \x/\lab in {1.5/{\frac12},3.2/{\frac32},4.9/{\frac52},6.6/{\frac72}}{
        % branch point
        \filldraw[black] (\x,0) circle (2.2pt);
        % cut to the right
        \draw[thick] (\x,0) -- (\x+0.85,0);
        % label
        \node[below=5pt] at (\x,0) {$u=\lab$};
    }
    % title-like annotation
    \node[align=left] at (6.6,1.8)
    {IR renormalons of pole mass};
\end{tikzpicture}
\caption{Pole mass 在 Borel \(u\)-平面上的 IR renormalon 解析结构示意图。黑点表示 branch points 的起点；从这些点向右延伸的线段表示相应的 branch cuts。}
\label{fig:pole_mass_borel_branch}
\end{figure}

具体地，定义Borel变换为$B[R](t)=\sum_{n=0}^{\infty} \frac{r_n}{n!}t^n$, 若$r_n=K a^n\Gamma(n+1+b)$,则
$$
B[R](t) \sim \frac{K \Gamma(1+b)}{(1-a t)^{1+b}} .
$$
若$B[R](t)$在正实轴没有奇点，且在无穷远处增长足够慢，则在形式上可定义Borel积分\footnote{这里的$\alpha$是微扰展开参数}，
\begin{equation}
R_{\rm Borel}(\alpha)=\int_{0}^{\infty}dt\,e^{-t/\alpha}B[R](t).
\end{equation}
如果 Borel 奇点落在正实轴上，那么 Borel 积分绕奇点上、下侧定义时会给出不同结果，从而产生歧义。这个歧义通常表现为
\begin{equation}
  \text{Im} R_{\rm Borel}=\mp(\pi K/a) e^{-1/(a\alpha)}(a\alpha)^{-b}.
\end{equation}
其中的正负号取决于积分路径位于正实轴的上侧还是下侧，极质量定义中的歧义性也反映在\textbf{Borel积分路径的歧义性}上，该路径积分的差异会给出一个非微扰的指数量级不确定性\footnote{对于式\eqref{eq: chatacter_Integral},a=2,b=0, 对于质量位移结果，部分经典文献中，将t重定义为u，即该积分在u=1/2处存在奇点。}。

设某物理量 $R(Q)$ 的 Borel 变换在正实轴上存在 IR renormalon 奇点 $u=u_0>0$。
则 Borel 积分的歧义一般具有形式
\[
\delta R \sim e^{-u_0/((-b_0)\alpha_s(Q))}.
\]
利用一圈 running coupling
\[
\alpha_s(Q)=\frac{1}{-b_0\ln(Q^2/\Lambda_{\rm QCD}^2)},
\]
可得
\[
e^{-u_0/((-b_0)\alpha_s(Q))}
=
e^{-u_0\ln(Q^2/\Lambda_{\rm QCD}^2)}
=
\left(\frac{\Lambda_{\rm QCD}^2}{Q^2}\right)^{u_0}.
\]
因此，位于 $u=u_0$ 的 IR renormalon 对应的非微扰歧义量级为
\[
\delta R(Q)\sim \left(\frac{\Lambda_{\rm QCD}^2}{Q^2}\right)^{u_0},
\]
再乘以 $R$ 自身的质量维数所需的硬标度因子。
对于 pole mass 而言，首个 IR renormalon 位于 $u_0=\frac12$，且
$m_{\rm pole}-m_{\overline{\rm MS}}$ 具有质量维数 1，
故
\[
\delta m_{\rm pole}\sim Q\left(\frac{\Lambda_{\rm QCD}^2}{Q^2}\right)^{1/2}
\sim \Lambda_{\rm QCD}.
\]
pole mass 定义上存在$\mathcal{O}(\Lambda_{\rm QCD})$的模糊性便来自leading阶的renormalon\footnote{对于 renormalon 奇点，需要区分奇点位置与奇点附近的局部结构。

\begin{itemize}
    \item 奇点位置 $u=u_0$ 由红外或紫外区域的幂次行为决定，而不由更高阶 running coupling 决定。
    一般地，若小动量区行为为$F(k^2)\sim k^a$,
    则对应的 IR renormalon 位置为
    $
    u_0=\frac{a}{2}.
    $
    \item 更高阶 RG 信息（如两圈、三圈 $\beta$ 函数）不会改变 $u_0$，而只会把 simple pole 改成 branch point，
    即把
    $
    \frac{1}{1-u/u_0}
    $
    修正为
    \[
    \frac{1}{(1-u/u_0)^{1+b}}
    \left[1+c_1(1-u/u_0)+c_2(1-u/u_0)^2+\cdots\right].
    \]
\end{itemize}}。

Subleading阶的renormalon会带来更高幂次的$\Lambda_{\rm QCD}$依赖，但对于pole mass而言，leading阶的$\mathcal{O}(\Lambda_{\rm QCD})$不确定性已经足以说明其定义上的模糊性问题\cite{Beneke:1998ui}。
\section{阈值质量方案：1S质量方案}
1S质量方案并非对极点质量的任意重命名，而是从非相对论 $Q \bar{Q}$ 阈值动力学中抽取出的短距质量。它定义为「虚构」$1^3 S_1$ 夸克偶素基态微扰质量的一半，因此天然吸收了总静能 $2 m_{\text {pole }}+V_{\text {stat }}$ 中的物理可控部分。由于 pole mass 的 $O\left(\Lambda_{Q C D}\right)$ 红外重整子与静态势中的对应长程修正在总静能中相互抵消， 1S质量方案也随之摆脱了极点质量方案中的领头红外重整子任意性。它适合描述 $q\bar{q}$ 等阈值过程；当将其用于单举衰变时，需要基于形式展开参数$\epsilon$进行来保证与极点质量方案中扰级数之间的红外敏感项逐阶一致地重组与对消。
\subsection{从单胶子交换到pNRQCD的静态势物理图像}


\subsection{1S质量方案}
    这里写如何规避红外重整子任意性问题，
\subsection{Upsilon Expansion}
A. H. Hoang等人在上世纪末本世纪初对于1S质量方案进行了系统的研究，提出了1S质量方案的定义和计算方法，并将其应用于顶夸克质量的衰变过程。通过修饰的微扰展开对$\alpha_s$进行重求和以消除固定阶的「红外重整子任意性问题」，具体地，将其用于inclusive decay之类的“硬过程”时，必须使用upsilon expansion，使非微扰束缚能级数与普通过程的微扰级数正确配对。
对于顶夸克的阈值对产生过程，1S质量方案提供了一个自然的框架来定义和计算顶夸克的质量。一个比较系统完整的质量转换关系被提出，并且在数值上展示了极好的收敛性(hep-ph/9904468(78))\cite{Hoang:1999zc}。

\begin{equation}
\begin{aligned}
M_{1 S}= & M_t-\epsilon \frac{M_t C_F^2 a_s^2}{8} \\
& -\epsilon^2 \frac{M_t C_F^2 a_s^2}{8}\left(\frac{a_s}{\pi}\right)\left[\beta_0(L+1)+\frac{a_1}{2}\right] \\
& -\epsilon^3 \frac{M_t C_F^2 a_s^2}{8}\left(\frac{a_s}{\pi}\right)^2\left[\beta_0^2\left(\frac{3}{4} L^2+L+\frac{\zeta_3}{2}+\frac{\pi^2}{24}+\frac{1}{4}\right)+\beta_0 \frac{a_1}{2}\left(\frac{3}{2} L+1\right)\right. \\
& \left.\quad+\frac{\beta_1}{4}(L+1)+\frac{a_1^2}{16}+\frac{a_2}{8}+\left(C_A-\frac{C_F}{48}\right) C_F \pi^2\right],
\end{aligned}
\label{eq:OS_to_1S_NNNLO}
\end{equation}
其中$L=\ln \left(\mu/\left(C_F a_s M_t\right)\right)$，$\epsilon$是upsilon expansion的形式展开参数，用于标记QCD微扰展开的阶数。具体地，微扰级数中的$\mathcal{O}(\alpha_s^N)$项被赋予$\epsilon^N$，而极质量到1S质量转换关系中的$\mathcal{O}(\alpha_s^N)$（除去对数项）则被视为$\epsilon^{N-1}$;按$\epsilon$重新展开并截断到固定阶后，最终再令该参数取1，在技术上以完成重求并实现消除红外重整子任意性问题的目的。
注意到，$M_{1S}/M_t$相对于1的偏离始于$\epsilon$的线性项，即$\mathcal{O}(\epsilon \alpha_s^2)$，此外，L的出现始于$\epsilon$的平方项。换而言之，$M_{t}=M_{1S}(1+\delta)$且$\delta \sim \mathcal{O}(\epsilon \alpha_s^2)$。代入L的表达式并展开至$\epsilon$的线性项，我们可以得到
\begin{equation}
    L=\ln \left(\mu/\left(C_F a_s M_t\right)\right)=\ln\left(\mu/\left(C_F a_s M_{1S}\right)\right)-\ln\left(1+\delta\right)=\ln\left(\mu/\left(C_F a_s M_{1S}\right)\right)-\delta+\frac{\delta^2}{2}+\mathcal{O}(\delta^3)
\end{equation}
若$L_{1S}$与$L_{OS}$的差异在$\mathcal{O}(\epsilon^n)$量级,则该差异出现于式\eqref{eq:OS_to_1S_NNNLO}中的$\epsilon^{n-1}$项,因此，$L_{1S}$与$L_{OS}$的差异对于$\mathcal{O}(\epsilon^{n-1})$的微扰级数展开截断而言可以忽略。

对于具体物理观测量中的质量方案的幂次分析的讨论，则需要根据具体的物理过程和观测量来进行分析。例如对于精确至两圈图水平微扰修正的D介子单举半轻衰变过程的衰变宽度如下：
\begin{equation}
\begin{aligned}
    \Gamma_{D_{i}} = \sum_{q = \mathrm{d}, \mathrm{s}} \hat{\Gamma}_0 \mrm{\left| V_{\mathrm{c}q} \right|^2} m_{\mathrm{c}}^5 \Big\{ 1 & + \frac{\alpha_s}{\pi} \frac{2}{3} \left( \frac{25}{4} - \pi^2 \right)\\
    & + \frac{\alpha_s^2}{\pi^2} \left[ \frac{\beta_0}{4} \frac{2}{3} \left( \frac{25}{4} - \pi^2 \right) \log \left( \frac{\mu^2}{m_{\mathrm{c}}^2} \right) + 2.14690 n_l - 29.88311 \right]\\  
    &- 8 \rho \delta_{\mathrm{s}q} - \frac{1}{2} \frac{\mu_{\pi}^{2}(D_{i})}{m_{\mathrm{c}}^2} - \frac{3}{2} \frac{\mu_{\mathrm{G}}^{2}(D_{i})}{m_{\mathrm{c}}^2} + 6 \frac{\rho_{\mathrm{D}}^3\left(D_{i}\right)}{m_{\mathrm{c}}^3} + \dots \Big\},
\end{aligned}
\label{eq:decaywidth-mass-scheme}
\end{equation}
我们以截断至NNNLO的贡献为例进行说明：对数项$\ln(\mu^2/m_{c}^2)$出现于$\alpha_s^2$项，则该对数项对于upsilon expansion展开级数的贡献始于$\mathcal{O}(\epsilon^2)$量级，因此方便起见可对该对数项单独进行质量方案的转换，仅需对质量转换公式精确至$\mathcal{O}(\delta^2)$，即保留至$\mathcal{O}(\epsilon \alpha_s^2)$量级即可\footnote{这里讨论对应的三个问题分别是：1. 能否把式\eqref{eq:OS_to_1S_NNNLO}
里L中的$M_{t}$直接换成$M_{1S}$; 2. 为什么代码中物理观测量里的log项没有直接应用质量转换公式而是保留了$1+\delta$的线性项之后再统一做$\epsilon$展开；3. 在1S质量方案下，NNNLO截断对应于什么样的丢弃规则。}。
\section{阈值减除质量方案: PS质量方案}
\section{标准UV短矩质量：$\bar{\text{MS}}$质量方案}
\section{不同质量方案的对比}
\newpage
\section{补充计算}
\subsection{不同重整化方案下质量位移的一般性结果}
本小节主要从重整化后的自能图出发，计算不同重整化方案下质量位移的一般结果，并进行简单讨论。对于重整化后的费米子两点函数，传播子的极点位置定义了所谓的 \emph{pole mass}。  
设重整化后的逆传播子为
\begin{equation}
S_R^{-1}(p)=\slashed p-m-\Sigma(\slashed p),
\end{equation}
其中 \(m\) 是某个给定重整化方案下的质量参数，\(\Sigma(\slashed p)\) 是重整化后的自能。

我们关心的问题是：

\begin{itemize}
    \item 如何由 \(\Sigma(\slashed p)\) 定义传播子的极点位置 \(m_{\rm pole}\)；
    \item 如何在一阶精度下导出
    \begin{equation}
    \delta m \equiv m_{\rm pole}-m
    =\frac{1}{4m}\text{Tr}[(\slashed p+m)\Sigma(\slashed p)]|_{p^2=m^2}
    +\mathcal{O}(\Sigma^2);
    \end{equation}
\end{itemize}
由洛伦兹协变性可知，在真空中的费米子二点函数只能依赖于 \(p^\mu\)，因此 \(\Sigma(\slashed p)\) 的最一般狄拉克结构只能由
$\mathbf{1},  \slashed p$两种\footnote{若理论本身宇称不守恒，且额外存在背景场，则需要重新讨论来自洛伦兹对称性的约束}。故可写为
\begin{equation}
\Sigma(\slashed p)=A(p^2)\,\slashed p+B(p^2)\,m,
\end{equation}
其中 \(A(p^2)\) 与 \(B(p^2)\) 都是无量纲函数。

传播子的极点位置由$\det S_R^{-1}(p)\Big|_{p^2=m_{\rm pole}^2}=0$定义。
等价地，也可以说在壳上存在旋量 \(u(p)\)，满足
\begin{equation}
S_R^{-1}(p)\,u(p)=0,
\qquad p^2=m_{\rm pole}^2.
\end{equation}

代入逆传播子的表达式，
\begin{equation}
\Big[
\bigl(1-A(p^2)\bigr)\slashed p
-
\bigl(1+B(p^2)\bigr)m
\Big]u(p)=0.
\end{equation}
在壳上满足\footnote{$S^{-1}\sim Z_2^{-1}(\slashed{p}-m_{\text{pole}} )+ \mathcal{O}(p^2-m_{\text{pole}}^2)$}
\begin{equation}
\slashed p\,u(p)=m_{\rm pole}\,u(p),
\end{equation}
因此可得到标量条件
\begin{equation}
\bigl(1-A(m_{\rm pole}^2)\bigr)m_{\rm pole}
-
\bigl(1+B(m_{\rm pole}^2)\bigr)m=0.
\end{equation}

定义
\begin{equation}
m_{\rm pole}=m+\delta m,
\qquad
\delta m=\mathcal{O}(\Sigma).
\end{equation}
由于 $A, B\sim \mathcal{O}(\Sigma)$ ，因此在一阶精度下，
\begin{equation}
A(m_{\rm pole}^2)=A(m^2)+\mathcal{O}(\Sigma^2),
\qquad
B(m_{\rm pole}^2)=B(m^2)+\mathcal{O}(\Sigma^2).
\end{equation}
因此标量条件在一阶精度下变成
\begin{equation}
\bigl(1-A(m^2)\bigr)(m+\delta m)-\bigl(1+B(m^2)\bigr)m=0
+\mathcal{O}(\Sigma^2).
\end{equation}
展开并整理后得到
\begin{equation}
\delta m
=
m\bigl[A(m^2)+B(m^2)\bigr]+\mathcal{O}(\Sigma^2).
\end{equation}
A和B是属于自能图的狄拉克结构的系数，可以凭此得到质量位移结果对自能图的依赖的一般性表达式。

使用正能投影算符$\slashed{p}+m$对自能图进行投影，并取迹，可以得到
\begin{equation}
\text{Tr}[(\slashed p+m)\Sigma(\slashed p)]|_{p^2=m^2}=4 p^2 A(m^2)+4 m^2 B(m^2)=4 m^2 \bigl[A(m^2)+B(m^2)\bigr].
\end{equation}
最终得到质量位移的一般计算结果
\begin{equation}
\boxed{
\delta m
=
\left.\frac{1}{4m}\text{Tr}\!\big[(\slashed p+m)\Sigma(\slashed p)\big]\right|_{p^2=m^2}
+\mathcal{O}(\Sigma^2)
}.
\label{eq: massdistance}
\end{equation}
该结果表明，质量位移 $\delta m$ 可以通过对自能图 $\Sigma(\slashed p)$ 进行特定的投影和取迹来计算\footnote{利用狄拉克代数的性质，使用求迹操作将质量位移相关的标量部分投影出来}, \textbf{式\eqref{eq: massdistance}的结构无重整化方案依赖}，因为它仅取决于：1. 自能图的狄拉克结构, 2. 极点条件: $S^{-1}u(p)=0$, 3. 自能修正的一阶展开。比如式\eqref{eq: massdistance}中的m是$\overline{\mathrm{MS}}$质量，$\Sigma$是$\overline{\mathrm{MS}}$方案下的自能图，那么式\eqref{eq: massdistance}就给出了$\overline{\mathrm{MS}}$质量与极点质量之间的关系。若m是极质量，则按照定义，$\delta m=0$，因此式\eqref{eq: massdistance}也给出了极质量方案下自能图的特定投影必须为零的条件。

传播子在极点附近应具有如下形式：
\begin{equation}
S_R(p)\xrightarrow[p^2\to m_{\rm pole}^2]{}
\frac{i\,Z_2^{\rm pole}}{\slashed p-m_{\rm pole}+i0}
+\text{regular}.
\end{equation}
这里 \(Z_2^{\rm pole}\) 是极点留数。若选择 OS 方案，则通常要求极点留数归一化为 \(1\)。这意味着：

\begin{itemize}
    \item 极点确实位于 \(p^2=m_{\rm pole}^2\)；
    \item 极点前的系数被规范为单位留数。
\end{itemize}

这两个条件共同决定了 OS 方案中的质量 counterterm 与波函数重整化常数。

从逆传播子的角度看，第一条是
\begin{equation}
S_R^{-1}(p)\,u(p)=0
\qquad \text{at }\slashed p=m_{\rm pole},
\end{equation}
而第二条则对应于在极点附近对 \(S_R^{-1}(p)\) 对 \(\slashed p\) 或 \(p^2\) 的一阶展开系数施加归一化条件。

在 OS（on-shell）方案中，重整化条件直接要求：

\begin{itemize}
    \item 重整化后的传播子极点位于物理质量 \(m_{\rm pole}\)；
    \item 极点留数为 \(1\)。
\end{itemize}

因此 OS 方案下的质量参数就是 pole mass：
\begin{equation}
m_{\rm OS}=m_{\rm pole}.
\end{equation}
这意味着自能中的相应修正被完全吸收到极点位置里。

但是在 QCD 中，pole mass 对软动量区域敏感，因此会带有长程红外敏感性\footnote{我们质量位移结果的红外结果进行量级分析时，需要在ekional 近似下进行讨论。此时式\eqref{eq: massdistance}的分子部分以及求迹操作仅会贡献到常数因子，不参与红外结构的分析}。具体地，
\begin{equation}
\delta m_{\mathrm{soft}} \sim \int \frac{d^d k}{(2 \pi)^d} \frac{1}{k^2} \frac{1}{v \cdot k} \sim \int k^3 \frac{1}{k^3} d k \sim \int_0^\mu d k \text { (constant) }\sim \mathcal{O}(\Lambda_{\rm QCD})
\end{equation}
\textbf{
一般地，重整化后的逆传播子可写为
\[
S_R^{-1}(\slashed p)=\slashed p-m-\Sigma(\slashed p).
\]
在 OS 重整化方案中，质量参数被定义为传播子的极点位置，因此
\[
S_{\rm OS}^{-1}(\slashed p)\big|_{\slashed p=m_{\rm pole}}=0,
\]
也就是说，自能修正对极点位置的贡献被完全吸收到 pole mass 的定义中，其中也包括软胶子的贡献。正因如此，pole mass 对红外区域敏感，从而具有$
\mathcal O(\Lambda_{\rm QCD})
$量级的本征不确定性。}